% Appendix B, from info/quartus_workflow.md: installing Quartus Prime Lite and taking a design from
% VHDL source onto the DE0-CV. One section per numbered section of the source, labelled q:N. The
% source is written for the instructor who demonstrates the board work; its opening is adapted
% here for a reader who has a board of their own.

\chapter{Quartus Prime Lite and the DE0-CV}
\label{app:quartus}

This appendix is the toolchain reference for the book's board work. \textbf{Nobody needs a board to
follow the book}: every exercise is checked on your own laptop with a testbench, and no chapter asks
you to run Quartus. In the course this book comes from, the Quartus and DE0-CV work is demonstrated
from the front, starting with the braking assistant in Chapter~\ref{c1:ch}, and no participant
needs Quartus or a board of their own. This appendix is for a reader who does have a DE0-CV and
wants to put the book's designs on it, the way the course demonstrates them. It lives here rather
than in a chapter because nobody following the book has to execute it.

It covers:
\begin{itemize}
  \item installing Quartus Prime Lite (\secref{q:2}).
  \item once installed, taking a VHDL file from source to a running FPGA:
  \begin{itemize}
    \item new project (\secref{q:3}).
    \item adding sources (\secref{q:4}).
    \item assigning pins (\secref{q:5}).
    \item compiling (\secref{q:6}).
    \item programming the board (\secref{q:7}).
  \end{itemize}
\end{itemize}

Everything described here ends on real FPGA hardware, verified by hand. That is a separate layer
from the exercises, which you verify on your own laptop by running a self-checking testbench under
GHDL (see \secref{c2:app:c}).

\section{What You Need}
\label{q:1}
\begin{itemize}
  \item \textbf{Quartus Prime Lite} (free edition): Intel/Altera's FPGA design software, covering:
  \begin{itemize}
    \item synthesis.
    \item placement and routing (``fitting'').
    \item programming file generation.
  \end{itemize}
  \item \textbf{Cyclone V device support}: a separate add-on package.
  \begin{itemize}
    \item The DE0-CV's FPGA, \code{5CEBA4F23C7N}, is a Cyclone V part.
    \item Quartus doesn't include every device family's support by default.
  \end{itemize}
  \item \textbf{USB-Blaster driver}: the DE0-CV programs over USB using Altera's USB-Blaster
    protocol.
  \begin{itemize}
    \item the driver usually installs alongside Quartus.
    \item it occasionally needs a manual nudge (\secref{q:8}).
  \end{itemize}
  \item A \textbf{Terasic DE0-CV board} and a USB cable.
\end{itemize}

Quartus Prime Lite is available for both Windows and Linux; the steps below are the same regardless
of host OS except where noted (the Windows 11-specific gotcha in \secref{q:2} doesn't apply on
Linux).

\section{Installing Quartus Prime Lite}
\label{q:2}
\begin{enumerate}
  \item Download Quartus Prime Lite from Intel's Quartus Prime Lite download page,
    \url{https://www.intel.com/content/www/us/en/software-kit/711791/intel-quartus-prime-lite-edition-design-software-version-20-1-for-windows.html}.
  \begin{itemize}
    \item That link is the Windows kit for one particular release. Any recent Lite release works,
      on either host OS; this book doesn't depend on a specific point version, so pick the edition
      matching your platform from the same download centre.
  \end{itemize}
  \item During installation, make sure \textbf{Cyclone V} is selected under the device family
    options.
  \begin{itemize}
    \item If Quartus is already installed without it, add Cyclone V support afterwards by either:
    \begin{itemize}
      \item re-running the installer.
      \item using Quartus's own ``Additional Software'' / device installer.
    \end{itemize}
  \end{itemize}
  \item On first connecting the DE0-CV board over USB, Windows should detect the USB-Blaster and
    either:
  \begin{itemize}
    \item install its driver automatically.
    \item prompt you to point it at Quartus's own driver files (typically under the Quartus
      installation directory).
    \item If Windows can't find a driver automatically, see \secref{q:8}.
  \end{itemize}
\end{enumerate}

\noindent\note[Windows 11 gotcha]a couple of Windows 11's security defaults can interfere with this
install or with programming the board later:
\begin{itemize}
  \item If the installer itself won't open:
  \begin{itemize}
    \item check Windows Security's \code{Smart App Control} (search for it in the Start menu).
    \item turn it off, then restart.
  \end{itemize}
  \item If everything installs cleanly but the board still won't program (\secref{q:7}):
  \begin{itemize}
    \item go to \code{Windows Security → Device security → Core isolation}.
    \item disable \textbf{Memory integrity}:
  \end{itemize}
\end{itemize}

\bookimage{0.62}{info/images/mem_integrity.png}{The \code{Memory integrity} option under Windows
11's Core isolation settings, switched off.}{q:fig:memintegrity}

\noindent Neither setting affects Quartus itself; both are general Windows 11 hardening defaults
that happen to be strict enough to interfere with driver-level USB access.

\section{Creating a New Project}
\label{q:3}
\begin{enumerate}
  \item Open Quartus and start the \textbf{New Project Wizard}
    (\code{File → New Project Wizard}).
  \item Choose a working directory and a project name.
  \begin{itemize}
    \item Quartus uses the project name as the default top-level entity name too.
    \item Keep them matching your VHDL entity's name to avoid having to override it later (for
      \code{or_gate.vhd}, name the project \code{or_gate}).
  \end{itemize}
  \item Choose \textbf{Empty Project} as the project type.
  \begin{itemize}
    \item This book doesn't use Quartus's IP catalog or any project templates.
  \end{itemize}
  \item Skip adding files at this step (\secref{q:4} covers it) unless you already know the exact
    file; either is fine.
  \item On the device selection page, select the exact part on the DE0-CV board:
  \begin{itemize}
    \item filter by family \textbf{Cyclone V}.
    \item then locate and select the exact device \code{5CEBA4F23C7N}.
  \end{itemize}
  \item Skip the EDA tool settings page (leave it on ``none''); this book runs its simulations
    directly in GHDL rather than through Quartus.
  \item Finish the wizard.
\end{enumerate}

\section{Adding Your VHDL Source}
\label{q:4}
\begin{itemize}
  \item Add your source file:
  \begin{itemize}
    \item \code{Project → Add/Remove Files in Project...}, then browse to your \code{.vhd} file and
      add it (for example, Chapter~\ref{c1:ch}'s \file{lectures/L01/or_gate/or_gate.vhd} or
      Chapter~\ref{c5:ch}'s \file{lectures/L05/parity_gen/parity_gen.vhd}). \code{or_gate} stands
      in for Chapter~\ref{c1:ch}'s braking assistant throughout this appendix because the
      assistant is built live in the course and deliberately not committed; the steps are
      identical either way.
  \end{itemize}
  \item If your design has more than one file (a top-level entity plus one or more
    subcomponents):
  \begin{itemize}
    \item add all of them.
    \item Quartus resolves the instantiation hierarchy automatically as long as every file is in
      the project.
  \end{itemize}
  \item Confirm the top-level entity is set correctly:
  \begin{itemize}
    \item \code{Project → Set as Top-Level Entity}, with your main file selected.
    \item For a single-file design like \code{or_gate} or \code{parity_gen}, this is automatic.
  \end{itemize}
\end{itemize}

\section{Assigning Pins with the Pin Planner}
\label{q:5}
Every port in your entity needs a physical pin on the FPGA before Quartus can produce a working
programming file, otherwise there's no way to connect your design's \code{a}/\code{b}/\code{x} (or
\code{bits}/\code{parity}) to an actual switch or LED on the board.
\begin{enumerate}
  \item Run \code{Processing → Start → Start Analysis & Elaboration} first, then open the Pin
    Planner:
  \begin{itemize}
    \item open it via \code{Assignments → Pin Planner}.
    \item on a brand-new project the port list is empty until Quartus has elaborated the design and
      knows its ports.
  \end{itemize}
  \item Each of your entity's ports appears as a row.
  \begin{itemize}
    \item for each one, fill in the \textbf{Location} column with the pin identifier for the
      physical switch or LED you want to use.
  \end{itemize}
  \item This book doesn't reproduce the DE0-CV's specific pin numbers here.
  \begin{itemize}
    \item Look them up yourself, either:
    \begin{itemize}
      \item directly in the Pin Planner (it lists the device's available pins).
      \item in the DE0-CV user manual's pinout tables.
    \end{itemize}
    \item Use whichever switches/LEDs are convenient on your board.
  \end{itemize}
  \item Set the \textbf{I/O Standard} column if Quartus doesn't infer a sensible default for the
    pin you chose.
  \begin{itemize}
    \item the DE0-CV's general-purpose I/O is 3.3\,V LVTTL/LVCMOS.
    \item consult the pin's entry in the Pin Planner if in doubt.
  \end{itemize}
  \item Save the assignments (they're written into the project's \code{.qsf} file).
  \begin{itemize}
    \item any change here requires recompiling (\secref{q:6}) before it takes effect.
    \item Quartus won't retroactively patch a \code{.sof} that was compiled before the pin
      assignment changed.
  \end{itemize}
\end{enumerate}

\section{Compiling the Design}
\label{q:6}
\begin{itemize}
  \item \code{Processing → Start Compilation} (or the $\blacktriangleright$ toolbar button) runs the
    full flow:
  \begin{itemize}
    \item analysis and synthesis.
    \item placement and routing (``fit'').
    \item timing analysis.
    \item assembly into a programming file.
  \end{itemize}
  \item Read the compilation report once it finishes, especially the \textbf{Warnings}:
  \begin{itemize}
    \item an unconnected port, an incomplete sensitivity list, or a latch inferred where you meant
      combinational logic all show up here as warnings rather than errors.
    \item they're easy to miss if you only check for a green checkmark.
  \end{itemize}
  \item A successful compile produces a \code{.sof} file (SRAM Object File) in the project's
    \code{output_files} directory.
  \begin{itemize}
    \item this is what \secref{q:7} loads onto the board.
    \item it's volatile: it configures the FPGA's SRAM directly and is lost on power-down; there's
      no persistent flashing step in this book.
  \end{itemize}
\end{itemize}

\section{Programming the DE0-CV Board}
\label{q:7}
\begin{enumerate}
  \item Connect the DE0-CV to your computer via USB and power it on.
  \item Open the Programmer: \code{Tools → Programmer}.
  \item Under \textbf{Hardware Setup}, select the USB-Blaster.
  \begin{itemize}
    \item it should be auto-detected if the driver installed correctly (see \secref{q:8} if not).
  \end{itemize}
  \item Make sure the \code{.sof} file produced in \secref{q:6} is listed, and its
    \textbf{Program/Configure} checkbox is ticked.
  \begin{itemize}
    \item if it's not listed, use \code{Auto Detect} or add the file manually.
  \end{itemize}
  \item Click \textbf{Start}.
  \begin{itemize}
    \item the progress bar should reach 100\% and report success.
  \end{itemize}
  \item Test the design directly on the board:
  \begin{itemize}
    \item flip the switches you assigned in \secref{q:5} and confirm the LED matches your entity's
      truth table:
    \begin{itemize}
      \item for \code{or_gate}: lit whenever either switch is on.
      \item for \code{parity_gen}: toggles every time exactly one switch changes.
    \end{itemize}
  \end{itemize}
\end{enumerate}

\section{Troubleshooting}
\label{q:8}
\begin{itemize}
  \item \textbf{USB-Blaster not detected in the Programmer}:
  \begin{itemize}
    \item check Windows Device Manager for an unknown device or one flagged with a driver error.
    \item point it manually at the USB-Blaster driver under the Quartus installation directory.
  \end{itemize}
  \item \textbf{Board won't program, no obvious error}:
  \begin{itemize}
    \item on Windows 11, disable \code{Memory integrity} under \code{Core isolation}
      (\secref{q:2}).
    \item this is the single most common cause once the driver itself is installed.
  \end{itemize}
  \item \textbf{Installer won't run or open}: check \code{Smart App Control} under Windows Security
    (\secref{q:2}).
  \item \textbf{Compilation succeeds but the design doesn't behave correctly on the board}:
  \begin{itemize}
    \item re-check the Pin Planner assignments (\secref{q:5}) before suspecting the VHDL.
    \item a swapped or missing pin assignment produces a board that compiles cleanly but behaves
      nonsensically, since it isn't wired the way you think it is.
  \end{itemize}
  \item \textbf{Fitter reports unassigned pins as errors}:
  \begin{itemize}
    \item some project settings treat this as fatal rather than a warning.
    \item either assign every port in the Pin Planner or, for a stray unused port during early
      development, mark it explicitly.
    \item Don't leave ports silently unassigned and expect Quartus to guess.
  \end{itemize}
\end{itemize}
